% Generated draft: requires mathematical and source review.
\documentclass[11pt]{article}
\usepackage[margin=25mm]{geometry}
\usepackage{fontspec}
\setmainfont{Latin Modern Roman}
\usepackage{amsmath,amssymb,mathtools,unicode-math}
\setmathfont{Latin Modern Math}
\usepackage{enumitem,needspace,xurl,hyperref}
\setlength{\parindent}{0pt}
\setlength{\parskip}{5pt}
\setlength{\emergencystretch}{3em}
\newcommand{\N}{\mathbb{N}}
\newcommand{\Z}{\mathbb{Z}}
\newcommand{\Q}{\mathbb{Q}}
\newcommand{\R}{\mathbb{R}}
\newcommand{\C}{\mathbb{C}}
\newcommand{\F}{\mathbb{F}}
\newcommand{\A}{\mathbb{A}}
\newcommand{\PP}{\mathbb{P}}
\newcommand{\E}{\mathbb{E}}
\newcommand{\eps}{\varepsilon}
\newcommand{\dd}{\,\mathrm{d}}
\begin{document}
\section*{UnsolvedMath ID: 800888}
\textbf{Source problem code:} catalogue entry
\subsection{Definitions and mathematical statement}
The source asks about a smooth totally real submanifold \(S\subset X\) of maximal real dimension \(n\). Here total reality means that \(X\) has complex dimension \(n\) and
\[
T_pS\cap J(T_pS)=\{0\}\qquad(p\in S),
\]
where \(J\) is the complex structure of \(X\).

The supplied excerpt does not specify the ambient hypotheses on \((X,\omega)\), the metric defining Hausdorff measure, or finiteness of that measure. A conditional formulation in the usual compact K\"ahler setting is as follows: let \((X,\omega)\) be a compact K\"ahler manifold of complex dimension \(n\), with \(\int_X\omega^n=1\), and fix the metric used to define \(\mathcal H^n\). Suppose \(S\subset X\) is a smooth totally real submanifold of real dimension \(n\), with \(0<\mathcal H^n(S)<\infty\). Define its normalized Hausdorff measure by
\[
\mu_S(A)=\frac{\mathcal H^n(A\cap S)}{\mathcal H^n(S)}
\]
for Borel sets \(A\subset X\).

An \(\omega\)-plurisubharmonic function is an upper semicontinuous, locally integrable function \(u\), locally equal to a plurisubharmonic function plus a smooth function, such that \(\omega+dd^cu\geq0\) as a current. Write \(\mathcal E(X,\omega)\) for the full Monge--Amp\`ere mass class, consisting of such functions with
\[
\int_X\left\langle(\omega+dd^cu)^n\right\rangle=\int_X\omega^n.
\]
The brackets denote the non-pluripolar Monge--Amp\`ere product; the normalization of \(dd^c\) must be fixed consistently.

The supplied source asserts that there is a unique \(\varphi_S\in\mathcal E(X,\omega)\) satisfying
\[
\left\langle(\omega+dd^c\varphi_S)^n\right\rangle=\mu_S,
\qquad \sup_X\varphi_S=0.
\]
\textbf{Question (Berman, as attributed in the source).} Determine the optimal regularity of \(\varphi_S\).

This remains an incompletely specified regularity question: the excerpt does not identify a regularity scale, distinguish global regularity from regularity on or away from \(S\), or specify compactness, closedness, or boundary conditions for \(S\). The compact K\"ahler and measure hypotheses above are explicit conditional reconstruction assumptions requiring confirmation from the source context. The source attributes non-pluripolarity of \(S\) to Sadullaev, reference [Sad], and presents existence and uniqueness as a consequence. Non-pluripolarity of a set means that it is not contained in the \(-\infty\) locus of a plurisubharmonic function locally, with the usual local interpretation on \(X\); this assertion alone does not state that its Hausdorff measure gives zero mass to every pluripolar subset. The needed measure-level justification is not supplied in the excerpt.
\subsection{Short English statement}
Determine the best regularity of the normalized \(\omega\)-plurisubharmonic potential whose Monge--Amp\`ere measure is normalized Hausdorff measure on a smooth, maximally totally real submanifold. The source leaves the precise ambient assumptions, boundary conditions, and intended regularity scale unresolved.
\subsection{Sources}
\begin{itemize}
\item Local UnsolvedMath registry, supplied complete source record: permanent ID \(800888\), exact source code \texttt{catalogue entry}, AMR Open Problem Lists. The registry labels the problem open; this is a supplied status, not an independently verified assessment.
\item Dinew, Guedj, and Zeriahi, \emph{Open Problems in Pluripotential Theory} (2015), Question 19. Supplied source URL: \url{https://arxiv.org/abs/1511.00705}. The record lists Dinew, Guedj, and Zeriahi as proposers and attributes the question/context to Berman.
\item Supplied provenance: extraction method \texttt{arxiv\_latex}; access date \texttt{2026-10-16}. These are supplied provenance claims; no source visit or independent verification is asserted here.
\item Sadullaev, reference [Sad] in the supplied excerpt, cited there for non-pluripolarity of \(S\). No title, publication details, or URL for [Sad] is supplied.
\end{itemize}
\end{document}
